\chapter{Integración}\label{ch:g12:integ}

La integración responde a dos preguntas a la vez: cuál es el área bajo una
curva y cómo se puede recuperar una \omterm{def:g11:func:function}{función} a partir de su tasa de
variación. El teorema fundamental del cálculo afirma que son la misma
pregunta: el descubrimiento más profundo y más útil de las matemáticas del
siglo XVII.

\section{La integral de una función continua}

\begin{definition}[Integral de una función positiva]\label{def:g12:integ:area}
Sea $f$ \omterm{def:g12:limcont:continuity}{continua} y positiva en $\intcc{a}{b}$. La
\emph{integral}\index{integral}
\[
\int_a^b f(x)\,\dd x
\]
es el área, en unidades de área, de la región limitada por la curva de $f$,
el eje $x$ y las rectas verticales $x = a$ y $x = b$.
\end{definition}

\begin{omfigure}
\begin{tikzpicture}
\begin{axis}[
  width=8.4cm, height=5cm,
  axis lines=middle, xmin=-0.4, xmax=3.4, ymin=-0.3, ymax=2.3,
  xtick={0.5,2.7}, xticklabels={$a$,$b$}, ytick=\empty,
  samples=100, domain=-0.2:3.3]
  \addplot[name path=curve, omDef, thick, domain=0.2:3.2]
    {1 + 0.5*sin(deg(2*x)) + 0.2*x};
  \path[name path=axis] (0.5,0) -- (2.7,0);
  \addplot[omDef!20] fill between[of=curve and axis, soft clip={domain=0.5:2.7}];
  \node at (axis cs:1.6,0.55) {$\displaystyle\int_a^b f$};
\end{axis}
\end{tikzpicture}
\end{omfigure}

Para una \omterm{def:g11:func:function}{función} de signo cualquiera, las áreas situadas por debajo del eje
$x$ se cuentan en negativo; y se conviene que
$\int_b^a f(x)\,\dd x = -\int_a^b f(x)\,\dd x$.

\begin{omfigure}
\begin{tikzpicture}
\begin{axis}[omaxis,
  width=9.6cm, height=4.8cm,
  xmin=-0.5, xmax=7, ymin=-1.4, ymax=1.4,
  xtick={3.1416,6.2832}, xticklabels={$\pi$,$2\pi$}, ytick={-1,1}]
  \addplot[name path=sine, omDef, thick, domain=0:6.2832, samples=120]
    {sin(deg(x))};
  \path[name path=xaxis] (0,0) -- (6.2832,0);
  \addplot[omDef!20] fill between[of=sine and xaxis,
    soft clip={domain=0:3.1416}];
  \addplot[omThm!20] fill between[of=sine and xaxis,
    soft clip={domain=3.1416:6.2832}];
  \node at (axis cs:1.57,0.4) {$+$};
  \node at (axis cs:4.71,-0.4) {$-$};
\end{axis}
\end{tikzpicture}

{\small Áreas con signo: $\int_0^{2\pi} \sin x\,\dd x = 0$, porque la región
que queda por debajo del eje (en rojo) cancela la que queda por encima (en
azul).}
\end{omfigure}

\begin{proposition}[Propiedades de la integral]\label{prop:g12:integ:properties}
Sean $f, g$ \omterm{def:g12:limcont:continuity}{continuas} en un \omterm{def:g10:numbers:interval}{intervalo} que contiene a $a, b, c$ y sea
$\lambda \in \R$.
\begin{enumerate}
  \item \emph{Linealidad:}
        $\displaystyle\int_a^b (f + \lambda g) = \int_a^b f + \lambda \int_a^b g$.
  \item \emph{Relación de Chasles:}
        $\displaystyle\int_a^c f = \int_a^b f + \int_b^c f$.
  \item \emph{Positividad:} si $a \leq b$ y $f \geq 0$ en $\intcc{a}{b}$,
        entonces $\displaystyle\int_a^b f \geq 0$; y si $f \leq g$,
        entonces $\displaystyle\int_a^b f \leq \int_a^b g$.
\end{enumerate}
\end{proposition}

\begin{proof}
La relación de Chasles y la positividad son inmediatas a partir de la
interpretación como área (las áreas se suman al yuxtaponer regiones; una
región de altura positiva tiene área positiva). La comparación se sigue
aplicando la positividad a $g - f$. La linealidad es intuitivamente clara
para la suma de funciones positivas (basta apilar las áreas) y se demuestra
con rigor en la universidad; \emph{aquí se admite}.
\end{proof}

\section{El teorema fundamental del cálculo}

\begin{definition}[Primitiva]\label{def:g12:integ:primitive}
Una \emph{primitiva}\index{primitiva} de $f$ en un \omterm{def:g10:numbers:interval}{intervalo} $I$ es una
\omterm{def:g11:func:function}{función} \omterm{def:g12:deriv:derivative}{derivable} $F$ en $I$ tal que $F' = f$.
\end{definition}

\begin{proposition}\label{prop:g12:integ:primunique}
Si $F$ es una \omterm{def:g12:integ:primitive}{primitiva} de $f$ en un \omterm{def:g10:numbers:interval}{intervalo} $I$, las \omterm{def:g12:integ:primitive}{primitivas} de $f$
en $I$ son exactamente las funciones $F + c$, con $c \in \R$. Dados
$x_0 \in I$ e $y_0 \in \R$, hay una única \omterm{def:g12:integ:primitive}{primitiva} con $F(x_0) = y_0$.
\end{proposition}

\begin{proof}
Si $G' = F' = f$, entonces $(G - F)' = 0$ en el \omterm{def:g10:numbers:interval}{intervalo} $I$, así que
$G - F$ es constante.\footnote{Que una \omterm{def:g11:func:function}{función} de \omterm{def:g12:deriv:derivative}{derivada} nula en un
\omterm{def:g10:numbers:interval}{intervalo} es constante se sigue del \cref{thm:g12:deriv:variations}: es a
la vez \omterm{def:g12:seq:monotonic}{creciente} y \omterm{def:g10:functions:variations}{decreciente}.} Recíprocamente, toda $F + c$ es una
\omterm{def:g12:integ:primitive}{primitiva}. La condición $F(x_0) = y_0$ fija la constante.
\end{proof}

\begin{theorem}[Teorema fundamental del cálculo]\label{thm:g12:integ:ftc}
\index{teorema fundamental del cálculo}
Sea $f$ \omterm{def:g12:limcont:continuity}{continua} en un \omterm{def:g10:numbers:interval}{intervalo} $I$ y sea $a \in I$. La \omterm{def:g11:func:function}{función}
\[
F \colon x \longmapsto \int_a^x f(t)\,\dd t
\]
es la \omterm{def:g12:integ:primitive}{primitiva} de $f$ en $I$ que se anula en $a$. En consecuencia, para
cualquier \omterm{def:g12:integ:primitive}{primitiva} $G$ de $f$ y para $a, b \in I$:
\[
\int_a^b f(t)\,\dd t = \bigl[G(t)\bigr]_a^b = G(b) - G(a).
\]
\end{theorem}

\begin{omfigure}
\begin{tikzpicture}
\begin{axis}[omaxis,
  width=9.2cm, height=5.6cm,
  xmin=-0.4, xmax=4.2, ymin=-0.4, ymax=2.9,
  xtick={0.5,2.2,2.7}, xticklabels={$a$,$x$,\ $x+h$}, ytick=\empty]
  \addplot[name path=curve, omDef, thick, domain=0:3.9, samples=80]
    {1 + 0.4*x + 0.2*sin(deg(2*x))};
  \path[name path=xaxis] (0,0) -- (3.9,0);
  \addplot[omDef!20] fill between[of=curve and xaxis,
    soft clip={domain=0.5:2.2}];
  \addplot[omProp!40] fill between[of=curve and xaxis,
    soft clip={domain=2.2:2.7}];
  \draw[black, dashed, thin] (axis cs:2.2,1.69) -- (axis cs:2.7,1.69);
  \draw[black, dashed, thin] (axis cs:2.2,1.925) -- (axis cs:2.7,1.925)
    node[above right, font=\small] {$f(x+h)$};
  \node at (axis cs:1.35,0.6) {$F(x)$};
\end{axis}
\end{tikzpicture}

{\small La idea de la demostración: el incremento $F(x+h) - F(x)$ es el
área de la banda estrecha (en naranja), encajada entre los rectángulos de
alturas $f(x)$ y $f(x+h)$ sobre una base de longitud $h$.}
\end{omfigure}

\begin{proof}[Demostración cuando $f$ es creciente]
Fijemos $x \in I$ y $h > 0$ con $x + h \in I$. Por Chasles,
\[
F(x+h) - F(x) = \int_x^{x+h} f(t)\,\dd t .
\]
Como $f$ es \omterm{def:g12:seq:monotonic}{creciente}, $f(x) \leq f(t) \leq f(x+h)$ para
$t \in \intcc{x}{x+h}$, y por comparación de \omterm{def:g12:integ:area}{integrales} (la \omterm{def:g12:integ:area}{integral} de una
constante $c$ sobre un \omterm{def:g10:numbers:interval}{intervalo} de longitud $h$ vale $ch$):
\[
h\,f(x) \leq F(x+h) - F(x) \leq h\,f(x+h),
\]
luego
\[
f(x) \leq \frac{F(x+h) - F(x)}{h} \leq f(x+h).
\]
Cuando $h \to 0^+$, $f(x+h) \to f(x)$ por \omterm{def:g12:limcont:continuity}{continuidad}, y el teorema del
sándwich da que el cociente incremental tiende a $f(x)$; el caso $h < 0$ es
simétrico. Por tanto, $F' = f$, y $F(a) = 0$. El caso \omterm{def:g12:limcont:continuity}{continuo} general (no
monótono) se demuestra en la universidad.

Por último, si $G$ es una \omterm{def:g12:integ:primitive}{primitiva} cualquiera, $G = F + c$
(\cref{prop:g12:integ:primunique}), luego
$G(b) - G(a) = F(b) - F(a) = \int_a^b f$.
\end{proof}

\begin{example}
$\displaystyle\int_0^1 x^2\,\dd x = \left[\frac{x^3}{3}\right]_0^1 =
\frac13$: el área bajo la \omterm{def:g11:quad:parabola}{parábola} es un tercio del cuadrado unidad, como
ya sabía Arquímedes.
\end{example}

La tabla de \omterm{def:g12:integ:primitive}{primitivas} se lee en la tabla de \omterm{def:g12:deriv:derivative}{derivadas} ($u$ designa una
\omterm{def:g11:func:function}{función} \omterm{def:g12:deriv:derivative}{derivable} y $c$, una constante arbitraria):
\begin{center}
\renewcommand{\arraystretch}{1.7}
\begin{tabular}{c|c||c|c}
$f(x)$ & \omterm{def:g12:integ:primitive}{primitiva} & $f$ & \omterm{def:g12:integ:primitive}{primitiva} \\
\hline
$x^n \ (n \neq -1)$ & $\dfrac{x^{n+1}}{n+1} + c$ &
$u'u^n \ (n \neq -1)$ & $\dfrac{u^{n+1}}{n+1} + c$\\
$\dfrac1x \ (x > 0)$ & $\ln x + c$ &
$\dfrac{u'}{u} \ (u > 0)$ & $\ln u + c$\\
$\eu^x$ & $\eu^x + c$ &
$u'\eu^u$ & $\eu^u + c$\\
$\cos x$ & $\sin x + c$ &
$\dfrac{u'}{\sqrt u} \ (u>0)$ & $2\sqrt u + c$\\
$\sin x$ & $-\cos x + c$ & &
\end{tabular}
\end{center}

\begin{method}[Reconocer la forma $u' \times (\text{algo en } u)$]\label{met:g12:integ:uprime}
Para integrar un producto, busca un factor que sea la \omterm{def:g12:deriv:derivative}{derivada} de una
\omterm{def:g11:func:function}{función} interior $u$, salvo una constante multiplicativa. Por ejemplo, en
$\int_0^1 x\,\eu^{x^2}\dd x$, el factor $x$ es $\frac12 (x^2)'$:
\[
\int_0^1 x\,\eu^{x^2}\dd x = \frac12\left[\eu^{x^2}\right]_0^1
= \frac{\eu - 1}{2}.
\]
\end{method}

\begin{theorem}[Integración por partes]\label{thm:g12:integ:ibp}
\index{integración por partes}
Sean $u, v$ \omterm{def:g12:deriv:derivative}{derivables} en $\intcc{a}{b}$ con \omterm{def:g12:deriv:derivative}{derivadas} \omterm{def:g12:limcont:continuity}{continuas}. Entonces
\[
\int_a^b u'(t)\,v(t)\,\dd t
= \bigl[u(t)\,v(t)\bigr]_a^b - \int_a^b u(t)\,v'(t)\,\dd t .
\]
\end{theorem}

\begin{proof}
La regla del producto da $(uv)' = u'v + uv'$; integrando los dos miembros
en $\intcc{a}{b}$ y usando el teorema fundamental en el primero se obtiene
$\bigl[uv\bigr]_a^b = \int_a^b u'v + \int_a^b uv'$.
\end{proof}

\begin{example}\label{ex:g12:integ:ibp}
$\displaystyle\int_0^1 t\,\eu^{t}\,\dd t$: tomamos $u' = \eu^t$ y $v = t$,
así que $u = \eu^t$ y $v' = 1$:
\[
\int_0^1 t\,\eu^t\,\dd t = \bigl[t\,\eu^t\bigr]_0^1 - \int_0^1 \eu^t\,\dd t
= \eu - (\eu - 1) = 1 .
\]
\end{example}

\section{Aplicaciones}

\begin{definition}[Valor medio]\label{def:g12:integ:mean}
El \emph{valor medio}\index{valor medio} de una \omterm{def:g11:func:function}{función} \omterm{def:g12:limcont:continuity}{continua} $f$ en
$\intcc{a}{b}$ ($a<b$) es
\[
\mu = \frac{1}{b-a}\int_a^b f(t)\,\dd t .
\]
\end{definition}

\begin{proposition}\label{prop:g12:integ:meanbounds}
Si $m \leq f \leq M$ en $\intcc{a}{b}$, entonces $m \leq \mu \leq M$.
\end{proposition}

\begin{proof}
Integramos las desigualdades $m \leq f(t) \leq M$ en $\intcc{a}{b}$ y
dividimos entre $b - a > 0$.
\end{proof}

\begin{method}[Área entre dos curvas]\label{met:g12:integ:areabetween}
Si $f \geq g$ en $\intcc{a}{b}$, el área entre las dos curvas es
$\int_a^b \bigl(f(x) - g(x)\bigr)\dd x$. Si las curvas se cortan, parte el
\omterm{def:g10:numbers:interval}{intervalo} por los puntos de corte e integra $\abs{f - g}$ trozo a trozo.
\end{method}

\begin{omfigure}
\begin{tikzpicture}
\begin{axis}[omaxis,
  width=8.4cm, height=5.8cm,
  xmin=-2.1, xmax=2.9, ymin=-1.8, ymax=5.6,
  xtick={-1,2}, ytick=\empty]
  \addplot[name path=line, omThm, thick, domain=-1.7:2.5] {x + 1};
  \addplot[name path=parab, omDef, thick, domain=-1.9:2.45] {x^2 - 1};
  \addplot[omDef!20] fill between[of=line and parab,
    soft clip={domain=-1:2}];
  \fill (axis cs:-1,0) circle (1.5pt);
  \fill (axis cs:2,3) circle (1.5pt);
\end{axis}
\end{tikzpicture}

{\small El área entre la recta $y = x + 1$ (en rojo) y la \omterm{def:g11:quad:parabola}{parábola}
$y = x^2 - 1$ (en azul) es
$\int_{-1}^{2}\bigl((x+1) - (x^2-1)\bigr)\dd x = \frac92$.}
\end{omfigure}

\section{Ejercicios}

\begin{exercise}[$\star$]\label{exo:g12:integ:1}
Calcula
\[
\int_1^2 \left(3x^2 - \frac{1}{x^2}\right)\dd x, \qquad
\int_0^{\pi/2} \cos t \,\dd t, \qquad
\int_0^{1} \frac{\dd t}{2t+1} .
\]
\end{exercise}

\begin{exercise}[$\star$]\label{exo:g12:integ:2}
Halla la \omterm{def:g12:integ:primitive}{primitiva} $F$ de $f(x) = x\eu^{x^2}$ en $\R$ tal que $F(0) = 1$.
\end{exercise}

\begin{exercise}[$\star$]\label{exo:g12:integ:3}
Calcula el \omterm{def:g12:integ:mean}{valor medio} de $f(t) = \sin t$ en $\intcc{0}{\pi}$ e interpreta
el resultado sobre una \omterm{def:g10:functions:graph}{gráfica}.
\end{exercise}

\begin{exercise}[$\star\star$]\label{exo:g12:integ:4}
Usando la \omterm{thm:g12:integ:ibp}{integración por partes}, calcula
\[
\int_1^{\eu} \ln t\,\dd t
\qquad\text{y}\qquad
\int_0^{\pi} t \sin t\,\dd t .
\]
\end{exercise}

\begin{exercise}[$\star\star$]\label{exo:g12:integ:5}
Calcula el área de la región comprendida entre la \omterm{def:g11:quad:parabola}{parábola} $y = x^2$ y la
recta $y = x + 2$.
\end{exercise}

\begin{exercise}[$\star\star$]\label{exo:g12:integ:6}
Sea $I = \displaystyle\int_0^1 \frac{\dd t}{1 + t}$.
\begin{enumerate}
  \item Calcula $I$.
  \item Para $n \in \N$, sea
        $I_n = \displaystyle\int_0^1 \frac{t^n}{1+t}\dd t$. Demuestra que
        $I_n + I_{n+1} = \dfrac{1}{n+1}$ y que
        $0 \leq I_n \leq \dfrac{1}{n+1}$.
  \item Deduce que
        $1 - \frac12 + \frac13 - \dots + \frac{(-1)^{n-1}}{n}
        \xrightarrow[n \to +\infty]{} \ln 2$.
\end{enumerate}
\end{exercise}

\begin{exercise}[$\star\star$]\label{exo:g12:integ:7}
La velocidad de un tren (en m/s) durante los primeros $100$ segundos tras
la salida se modeliza mediante $v(t) = 30\bigl(1 - \eu^{-t/50}\bigr)$.
Calcula la distancia recorrida durante esos $100$ segundos y la velocidad
\omterm{def:g11:stat:mean}{media} del tren en ese \omterm{def:g10:numbers:interval}{intervalo}.
\end{exercise}

\begin{exercise}[$\star\star\star$]\label{exo:g12:integ:8}
Para $n \geq 1$, sea
$S_n = \dfrac1n \displaystyle\sum_{k=1}^{n} \frac{1}{1 + k/n}$.
\begin{enumerate}
  \item Interpreta $S_n$ como un área de rectángulos que aproxima una
        región bajo la curva de $t \mapsto \frac{1}{1+t}$ en
        $\intcc{0}{1}$.
  \item Usando la monotonía de $t \mapsto \frac{1}{1+t}$, demuestra que
        \[
        S_n \leq \int_0^1 \frac{\dd t}{1+t} \leq S_n + \frac1n\left(1 - \frac12\right),
        \]
        y deduce que $\lim\limits_{n\to+\infty} S_n = \ln 2$.
\end{enumerate}
\end{exercise}

\begin{exercise}[$\star\star\star$]\label{exo:g12:integ:9}
\emph{(\omterm{def:g12:integ:area}{Integrales} de Wallis.)} Para $n \in \N$, sea
$W_n = \displaystyle\int_0^{\pi/2} \sin^n t\,\dd t$.
\begin{enumerate}
  \item Calcula $W_0$ y $W_1$.
  \item Escribiendo $\sin^{n+2}t = \sin t \cdot \sin^{n+1} t$ e integrando
        por partes, demuestra que $W_{n+2} = \dfrac{n+1}{n+2}\,W_n$.
  \item Deduce $W_2$, $W_3$ y $W_4$, y demuestra que $(W_n)$ es \omterm{def:g10:functions:variations}{decreciente}
        y positiva.
\end{enumerate}
\end{exercise}

\section{Problema: Arquímedes contra la máquina}

\begin{problem}[{Problema de fin de semana --- el área bajo la parábola,
calculada de tres maneras a lo largo de veintidós siglos, y la identidad
secreta del logaritmo como área}]
\label{pb:g12:integ:1}
Hacia el año 240 a.\,C., Arquímedes calculó el área exacta de un segmento
parabólico sin \omterm{def:g10:coordgeom:system}{coordenadas}, sin límites y sin álgebra. Diecinueve siglos
después, los rectángulos de Riemann lo rehicieron a base de encajonarlo; y
el teorema fundamental del cálculo (\cref{thm:g12:integ:ftc}) lo hace hoy
en una línea. Este problema juega los tres partidos y usa después la misma
máquina para revelar lo que es realmente el \omterm{def:g12:exp:ln}{logaritmo}: un área con una
simetría de escala.

\medskip
\noindent\textbf{Parte I --- Soltura.}
\begin{enumerate}
  \item Calcula $\displaystyle\int_0^1 (3x^2 - 2x + 1)\,\dd x$, \
        $\displaystyle\int_1^{\eu} \frac{\dd x}{x}$ \ y
        $\displaystyle\int_0^{\pi/2} \cos x\,\dd x$.
  \item Localiza las formas $u'u$ (\cref{met:g12:integ:uprime}):
        $\displaystyle\int_0^1 x\,\eu^{x^2}\dd x$ y
        $\displaystyle\int_0^1 \frac{2x}{x^2 + 1}\,\dd x$.
  \item Por partes (\cref{thm:g12:integ:ibp}):
        $\displaystyle\int_0^1 x\,\eu^x \dd x$ y
        $\displaystyle\int_1^{\eu} \ln x\,\dd x$.
  \item Calcula el \omterm{def:g12:integ:mean}{valor medio} (\cref{def:g12:integ:mean}) de $\sin$ en
        $\intcc{0}{\pi}$ y observa que \emph{no} es $\frac12$.
  \item Calcula el área entre la recta $y = x$ y la \omterm{def:g11:quad:parabola}{parábola} $y = x^2$ en
        $\intcc{0}{1}$ (\cref{met:g12:integ:areabetween}).
\end{enumerate}

\medskip
\noindent\textbf{Parte II --- La \omterm{def:g11:quad:parabola}{parábola}, de tres maneras.}
\begin{enumerate}[resume]
  \item El montaje de Riemann para el área bajo $y = x^2$ en
        $\intcc{0}{1}$: escribe la suma inferior $L_n$ y la suma superior
        $U_n$ sobre $n$ rectángulos iguales (como en el
        \cref{exo:g12:integ:8}).
  \item Demuestra por inducción (con la máquina del \cref{pb:g12:seq:1}) la
        fórmula de la suma de cuadrados
        \[
        1^2 + 2^2 + \dots + n^2 = \frac{n(n+1)(2n+1)}{6} .
        \]
  \item Deduce formas cerradas para $U_n$ y $L_n$, calcula su límite común
        y concluye: el área es $\frac13$.
  \item Ahora la máquina: calcula $\int_0^1 x^2\,\dd x$ con el teorema
        fundamental, en una línea. Compara los esfuerzos.
  \item Arquímedes lo enunciaba de otra manera: \emph{un segmento
        parabólico es $\frac43$ del triángulo inscrito}. Para el segmento
        que recorta en $y = x^2$ la cuerda que une $(-1, 1)$ con $(1, 1)$:
        calcula el área del segmento con una \omterm{def:g12:integ:area}{integral}, el área del
        triángulo inscrito (con vértice en $(0, 0)$) y comprueba la \omterm{def:g11:seq:geometric}{razón}
        del maestro.
  \item El método del propio Arquímedes: rellena el segmento con el
        triángulo grande $T$, después con dos triángulos que suman
        $\frac T4$, después con cuatro que suman $\frac{T}{16}$, y así
        sucesivamente. Suma la serie geométrica y recupera $\frac43 T$: la
        tableta de chocolate mordida sin fin del volumen anterior,
        devorada por un geómetra griego.
  \item Una frase para cada uno: la exhausción (Arquímedes), el
        encajonamiento (Riemann) y la búsqueda de \omterm{def:g12:integ:primitive}{primitivas}
        (Newton--Leibniz). ¿Qué necesita cada método y qué entrega?
\end{enumerate}

\medskip
\noindent\textbf{Parte III --- El \omterm{def:g12:exp:ln}{logaritmo} es un área.}
Para $x > 0$, tomemos $A(x) = \displaystyle\int_1^x \frac{\dd t}{t}$.
\begin{enumerate}[resume]
  \item Da $A(1)$ y $A'(x)$ (\cref{thm:g12:integ:ftc}) y concluye que $A$
        es exactamente el \omterm{def:g12:exp:ln}{logaritmo neperiano} de la
        \cref{def:g12:exp:ln}.
  \item El milagro de la escala: fija $a > 0$ y estudia
        $g(x) = A(ax) - A(x)$. Calcula $g'$ (regla de la cadena), deduce
        que $g$ es constante, evalúa la constante y concluye la \omterm{def:g10:algebra:equation}{ecuación}
        funcional
        \[
        \ln(ab) = \ln a + \ln b ,
        \]
        demostrada por puro cálculo: el área de $1$ a $ab$ se parte en
        copias reescaladas.
  \item Deduce de la pregunta 14: $\ln(a^n) = n\ln a$ y
        $\ln\frac1a = -\ln a$.
  \item El \cref{exo:g12:integ:8} lee
        $\frac{1}{n+1} + \frac{1}{n+2} + \dots + \frac{1}{2n}$ como
        rectángulos bajo $\frac{1}{1 + t}$: calcula esa suma para $n = 10$
        (con tres decimales) y compárala con $\ln 2$. ¿Qué sumas de sabor
        armónico, divergentes término a término, convergen aquí hacia un
        área?
\end{enumerate}

\medskip
\noindent\textbf{Parte IV --- Acumulación.}
\begin{enumerate}[resume]
  \item Un coche acelera con velocidad $v(t) = 3t^2$ m/s para
        $t \in \intcc{0}{10}$. Calcula la distancia recorrida, la velocidad
        \omterm{def:g11:stat:mean}{media} y el instante en que la velocidad instantánea es igual a la
        \omterm{def:g11:stat:mean}{media}.
  \item Una varilla de $4$ metros tiene densidad lineal
        $\rho(x) = 2 + x$ kg/m. Calcula su masa total y su centro de masas
        $\dfrac{\int_0^4 x\,\rho(x)\,\dd x}{\int_0^4 \rho(x)\,\dd x}$: el
        punto de equilibrio del triángulo de cartón del volumen anterior,
        calculado por fin con pesos que varían.
  \item \omterm{def:g12:integ:area}{Integrales} de Wallis (\cref{exo:g12:integ:9}): calcula $W_0$, $W_1$
        y $W_2 = \int_0^{\pi/2} \sin^2 t\,\dd t$ usando la linealización
        del \cref{pb:g12:trigo:1}. (Su escalera infinita sube, en los
        volúmenes universitarios, hasta una fórmula en producto para $\pi$
        y hasta la normalización de la campana de Gauss.)
  \item Final: las tres caras de la integración: un \emph{área} por
        definición, una \emph{acumulación} en el mundo (distancia, masa) y
        una \emph{\omterm{def:g12:integ:primitive}{primitiva}} por el teorema fundamental; y su historia en
        tres nombres. Cierra con la perla del capítulo: ¿qué tecla
        cotidiana de la calculadora es en secreto el área bajo $\frac1t$ y
        qué identidad demostró esa área?
\end{enumerate}
\end{problem}
